\honest{Changing the Definition of Science}{Eliezer Yudkowsky}{2008}

I quote New Scientist. Max Tegmark says the multiverse stands or falls with the tested theories it follows from. The philosopher Colin Howson wants to drop Popper's deductive account of science for one based on degrees of belief, with ``a fundamental connection between the subjective concept of belief and the cold, hard mathematics of probability.''

``I'm a good deal less of a lonely iconoclast than I seem. Maybe it's just the way I talk.'' \nb{A Bayesian philosophy of science was well established: Howson and Urbach's book first appeared in 1989. This is the first post in the sequence to point to that literature.}

My differences from the mainstream Bayesians are three. \nb{The article's own phrase, ``the subjective concept of belief,'' names a fourth. The Stanford Encyclopedia cites Howson for the view that Bayesians need not say what priors we should have, except that they be coherent. Two days earlier, in ``Science Isn't Strict Enough,'' Yudkowsky had written that not knowing the priors ``doesn't mean you get a free, personal choice of making the priors whatever you want.''}

First, I am not in academia, so I can say ``extreme'' things that others ``might well already be thinking.''

Second, ``just teaching probability theory won't be nearly enough.'' We must combine it with the study of cognitive biases and social psychology into a new ``Art of Bayescraft.'' My evidence: Robert Aumann, who first proved that Bayesians with the same priors cannot agree to disagree, is a believing Orthodox Jew. \nb{Aumann describes religion as ``an experience,'' ``mainly an emotional and aesthetic one,'' and ``orthogonal'' to science. The post assumes it is a failure of probabilistic reasoning.} This ``is not something I've seen suggested elsewhere.'' \nb{By the 1980s, research on judgment and decision making had paired normative models such as probability theory with descriptive models of how people err, and with prescriptive ``designs for improvement,'' which came to include teaching meant to counteract biases.}

Third, it is possible to do far better than science. A Bayesian superintelligence, seeing its first apple fall, could invent calculus, generalize Newton's laws, notice their action at a distance, and ``consider that General Relativity might be worth testing.'' I give no argument for this here. \nb{Yudkowsky argues it four days later, in ``That Alien Message.''} Our minds are so noisy that they need ``orders of magnitude more extra evidence to set us back on track after we derail,'' and academia ``is even slower.''
